\documentclass[11pt]{article}
\usepackage[margin=1in]{geometry}
\usepackage{microtype}
\usepackage{amsmath}
\usepackage[numbers,sort&compress]{natbib}
\usepackage{hyperref}
\usepackage{parskip}

\bibliographystyle{unsrtnat}

\begin{document}

\section*{Background and Central Hypothesis}

We conceptualize intelligent agents as solving problems through a sequence of two steps. When faced with a concrete problem (like AlphaZero with a board state, or a language model with a math problem), the agent has an initial set of priors that inform the next action it takes (the next move to be played or the next token to be generated). The agent then uses an exploratory phase to accumulate information and shifts its prior distribution over the next action towards something (hopefully) better. Formally, this is the structure of Bayes-adaptive sequential decision-making \citep{duff2002,ghavamzadeh2015}.

In AlphaZero \citep{silver2018alphazero}, this exploratory phase is hardcoded through a Monte-Carlo tree search of the game tree where each position is scored using the learned value function. Once this exploration is finished, there is also a hardcoded algorithm to use this valued game tree to improve its initial prior distribution over the next move.

In language models with ``reasoning'' \citep{deepseek2025r1}, this exploratory phase is simply the chain-of-thought the model uses, and is determined purely by the model itself. Once this phase is completed, the model leverages its in-context learning (ICL) abilities \citep{brown2020gpt3,xie2022icl} to improve its answer. The key difference from AlphaZero is that both how the exploration is executed, and how the exploration is used to improve the prior distribution over actions, is learned during training, instead of being a hardcoded algorithm. This difference informs our central hypothesis.

\paragraph{Central hypothesis.} When training jointly shapes both \emph{(i)} how an agent explores and \emph{(ii)} how it updates its policy from the exploration --- as in RL-trained reasoning models, but not in AlphaZero --- gradient pressure incentivises two internal representations to co-develop: an explicit representation of the agent's beliefs over task-relevant latents, and a causal model of the environment in which the agent itself figures as a distinct, well-separated cause. 

\paragraph{Why we expect this.} Optimal exploration is exploration that maximally increases information about the agent's task-relevant beliefs \citep{russo2014ids}. To do this, the agent must choose exploratory actions informed by its current beliefs, which motivates an internal representation of those beliefs and circuits that consult them to guide exploration. The second representation is required for a related reason: when the agent's own actions are part of what produces its next observation (as is the case for every token of a chain-of-thought), correctly attributing the outcome to its real causes requires the agent to distinguish what it produced itself from what is due to the world. Without this distinction, gradient updates conflate self-generated patterns with genuine signal about the task. When both the exploration and the update from it are themselves learned, gradient descent has every reason to develop both representations.

The first half of this hypothesis regarding internal belief representations has precedent in two strands of prior work. In the setting of sequence prediction in the context of Hidden Markov Models (HMMs), \citet{shai2024belief} show that transformers trained as next-token predictors develop linearly recoverable internal representations of the distribution over the hidden states. In active settings, \citet{hennig2023emergence} show that RL-trained recurrent networks develop similar representations over hidden states, and \citet{mikulik2020bayesopt,ortega2019meta} argue, on behavioural and representational grounds, that meta-trained agents converge to Bayes-optimal solutions whose state must encode a posterior over the task.

The second half of the hypothesis regarding a model of the self as a causal agent is, to our knowledge, new as a mechanistic claim. Its motivation is that when the agent's own action is one of the causes of its next observation, as is the case for every token of a chain-of-thought, correctly attributing the informational content of that observation requires distinguishing self-caused from world-caused signal. Initial evidence that post-trained language models do encode such self-representations comes from \citet{betley2025tellme} on behavioural self-awareness, \citet{panickssery2024self} on self-recognition correlating with self-preference, and the forthcoming work of one of us (Asvin) with Jack Lindsey, in which we find that post-trained models can detect when they are reading their own on-policy generations and modulate behaviour accordingly. 

The contribution of the proposed work --- spanning both halves of the hypothesis --- is to test mechanistically, through circuit-level probes and interventions, whether the belief representations of \citet{shai2024belief} and \citet{hennig2023emergence} are entangled with the self-representations of \citet{betley2025tellme}, whether the two factorise as a unified Self/World structure, and whether this structure emerges from the joint training pressure on exploration and update that distinguishes RL-trained reasoning models from AlphaZero-style systems. We pursue this first in a controlled toy setting where the analytical ground truth is computable, and then extend the same probe suite to open-source language models trained with environment interaction. The combination of circuit-level interventions, systematic variation of training regime, and joint examination of belief and self representations is, to our knowledge, what is new: existing work treats these elements separately, studies single regimes, and remains largely behavioural or decoding-based.

Both halves of the central hypothesis have well-studied developmental analogues, which we take as motivation. The first half --- the emergence of belief representations under exploration pressure --- has a substantial behavioural literature: infants attend preferentially to stimuli at intermediate complexity \citep[the Goldilocks effect,][]{kidd2012goldilocks}, select informative causal interventions \citep{lapidow2020informative}, track their own learning progress \citep{poli2020infants}, adapt their search to the ecological structure of the environment \citep{ruggeri2022ecological}, and show pupillometric signatures of active prediction of upcoming stimuli \citep{zhang2019prediction}. What developmental work cannot resolve, however, is the internal representations underlying these behaviours: infants behave as if they are optimising for information, but the objectives themselves and the internal states that implement them are not directly accessible from outside.

The second half --- the self / world distinction --- has equally direct developmental precedents, through a mechanism that maps especially cleanly onto the one we appeal to in language models. Within hours of birth, infants alter their sucking rhythm to produce preferred sounds \citep{decasper2009lateralized}. Newborns will raise an arm into a light beam and attend to the resulting hand shadow \citep{vandermeer1997limelight}. By three months, infants kick more vigorously when their ankle is tethered to a mobile that responds to their movement than under a yoked non-contingent control \citep{rovee1969conjugate}, and generalise the same principle across modalities and effectors \citep{lee2013contingent}. The mechanism running through all of these is temporal contingency detection: the infant learns to identify its own actions by their reliable, short-latency coupling to reafferent feedback. This is precisely the causal-sampling property we use as the signature of self in our toy and LLM analyses --- a self-emitted token is recognised as self by being a sample from the model's own prior output distribution, in the same way that an infant's movement is recognised as self by its temporal coupling to the motor intention that preceded it.

More striking still is the prenatal evidence, an argument developed at greater length by one of us \citep{rutar2024prenatal}. From the onset of foetal movement around 7.5 weeks gestation \citep{prechtl1985ultrasound}, the structure of motor output increases rapidly: early vermicular movements give way to fractionated, spatially selective limb and joint activity by 20 weeks \citep{devries1982emergence,piontelli2010conclusions,tiriac2014selfgenerated,lindemann2016spindle}, consistent with a shift from coarse, metabolically cheap sampling toward more targeted exploration of the foetus's own sensorimotor space. The earliest behavioural evidence of a primitive self-model appears at approximately the same time: by 22 weeks foetal hand trajectories toward the mouth follow more direct paths than those toward the eyes \citep{zoia2007evidence}, implying a mapping from motor commands to their expected sensory consequences on a self-structured target. The co-emergence of structured exploration and a primitive causal model of the body is consistent with the pattern our hypothesis predicts, though as a correlational observation it cannot show that one drives the other. This is, in the end, where developmental data runs out: it can document the behavioural signatures of the trajectory we hypothesise, but it cannot vary the training regime, intervene on representations, or access the ground truth. We see the flow of influence as bidirectional. Developmental research supplies the empirical patterns and conceptual vocabulary that motivate our hypothesis; in turn, by investigating in a system where the ground truth is computable and the training objective is under experimental control, our work aims to clarify which causal mechanisms can produce the developmental patterns and which cannot, sharpening the hypotheses the developmental community is itself testing.

\section*{Experimental Design for the First Half of the Year}

In order to isolate the various factors above cleanly, we propose a first experiment where we modify the setup in \citet{shai2024belief}. They train a transformer network on a next-token prediction task, where the tokens are generated by a Hidden Markov Model (HMM). They find that purely through being trained on the next prediction task, the model learns to internally represent a distribution over the hidden states of the HMM, and to learn correct update rules over this distribution from one token to the next. This is the simplest model of a \emph{pure predictor}. The model cannot influence the token stream in any way, and is constrained to passive prediction. As such, there is no "self" or "agency" to be learned. 

We propose to modify the set-up minimally in order to introduce agency. We maintain the same HMM model as before for the world, but introduce an ability for the model to "act". The token stream is now of the form $(a_1,o_1,\dots,a_n,o_n)$ where the $a_i$ are "actions", i.e., they are tokens sampled from the prediction of the neural network on the task. These actions \emph{cause} a transition in the hidden state of the HMM, and optionally provide a reward to the agent. The $o_i$ are observations as before and emitted when the HMM undergoes a state transition.

The generative process underlying this token stream factorises naturally into two coupled Markov processes. A \emph{World} chain governs the hidden environment state, with transitions driven by the agent's actions and emissions giving the observations $o_i$. A \emph{Self} chain governs the agent's own policy state, with transitions driven by the observations it receives and emissions giving the actions $a_i$. The joint hidden state is the product $\text{Self} \times \text{World}$, with token-position parity determining whose turn it is. Structures of this kind appear in the literature under different names (the World half is an input-output HMM \citep{bengio1995iohmm} with actions as inputs, and the joint system is a coupled HMM \citep{brand1996coupled}) though the explicit Self / World decomposition is not standard. 

We refer to this joint generative model as the \textbf{SWITCH} (Self/World Interactive Turn-taking Coupled HMM). We make it explicit here because the central hypothesis predicts that the trained transformer should recover this factorisation \emph{internally}: a Self subspace whose activations encode the agent's policy state, a separable World subspace whose activations encode the Bayes-adaptive posterior over the hidden environment state, and a clean separation between the two. Recovery of this factorisation would simultaneously validate both halves of the hypothesis and tie our result to the broader Simplex line of work on internal factorisation of independent generative processes.

This setup has several advantages for the purpose of testing our hypothesis. First, the notion of an "action" needs to be learned by the model - it is not informed ahead of time that the $a_i$ are causally downstream of its outputs. Second, we can vary the training regime to see whether the internal representations change under the meta-learning pressures we discussed above. We propose five different training set-ups.

\begin{enumerate}
    \item \emph{Pure prediction.} We train the model on token sequences emitted by the mathematically optimal agent in this set-up. This can be mathematically computed, as in \citet{shai2024belief}. In this case, we dissociate the action tokens from the outputs of the model, and we expect the model to learn no representations of self or agency.
    \item \emph{Predictive control.} The model produces actions autoregressively during training, so the action stream entering the input is causally produced by the model. We will run two sub-regimes that differ only in which token positions contribute to the cross-entropy loss: in (2a) the loss is applied at both action and observation positions, so the model is trained to predict both its own next action and the next observation the world would emit; in (2b) the loss is applied only at observation positions, so the action positions contribute to the input but not to the gradient, and the model is trained purely as a world model conditioned on its own action history. The two sub-regimes ask whether the structure of the action-observation interaction alone is enough to induce a self-representation (2b), or whether the model additionally needs an explicit gradient at the action positions (2a). We hypothesize that the self-representation emerges in both, but that the gradient at action positions in (2a) leads to a sharper and more robust representation.
    \item \emph{Trajectory training through RL.} We train the model through reinforcement learning in order to maximize reward over long trajectories. In this case, we don't directly provide a signal for the model about its token predictions, but nevertheless expect the model to learn an internal representation of the SWITCH.
    \item \emph{Active exploration.} We allow the model to ``explore'' before being scored on its next-token prediction. The concrete instance we will run first is a fixed exploration budget: at each prediction step the model takes $K$ actions and receives the corresponding $K$ observations before being scored on the prediction of the next observation. We vary $K$ from $1$ (essentially Regime 2b) upwards, and ask how the internal representations and the prediction quality depend on the exploration budget. Other instances --- variable budgets, model-decided stopping, exploration combined with an explicit reward --- are natural extensions once the fixed-$K$ version is understood.
    \item \emph{Mixed actors.} The active regimes so far always have the model itself as the actor at every action position. We now expose the model to action streams produced by a \emph{family} of policies, one of which is itself. Each trajectory is generated by first sampling an actor identity (the model itself, a random policy, the optimal Bayes-adaptive policy, an earlier checkpoint, or another deliberately distinct policy) and then having that actor produce all action tokens. The model is not told which actor produced any given trajectory and is trained on observation-prediction loss as in (2b). To predict observations well, the model must infer something about the actor's policy from the action statistics, and the natural feature is an \emph{actor-identity} representation that distinguishes the model's own action stream from those of other actors. This is to be contrasted with the \emph{token-type} self-representation (action versus observation positions) that previous regimes already encourage. This regime is also the toy analogue of the LLM training story, where the model is a predictor over many voices in pretraining and becomes the privileged producer of one of them in post-training. Crucially, by varying the actor while holding the World fixed, Regime 5 is the cleanest test of the causal form of the second hypothesis prediction: it forces the model to learn a World invariant under actor identity --- the structural signature of a genuine causal separation between self and world, distinct from the action-position / observation-position asymmetry which the other regimes already provide.
\end{enumerate}

\section*{Analysis of Trained Networks}

For each of the five training regimes above we obtain a trained transformer whose internal computations we can interrogate. Because the SWITCH is small, the exact Bayes-adaptive posterior over its World hidden state at every timestep is analytically computable; we therefore have a ground-truth comparator that experiments on natural-language reasoning models do not have. Our analyses test the two claims of the central hypothesis separately: (i) that the model develops an internal representation of its beliefs over the World hidden state of the SWITCH, and (ii) that the model develops a representation of itself as a distinct cause of the trajectory it observes. We then ask, across regimes, whether these representations emerge under the training pressures the hypothesis identifies.

\paragraph{Belief geometry.} Following \citet{shai2024belief}, we propose training linear probes from residual-stream activations at each layer and position to the exact Bayes-adaptive posterior over the World hidden state given history. The first question is whether the mixed-state-presentation geometry recovered by \citet{shai2024belief} in the passive case generalises to the active case at all. The second is whether the posterior the model encodes is the right one: in the active regimes the Bayes-adaptive posterior must condition on the action that was taken, which is itself a sample from the model's own output distribution. If the model performs this update correctly, it must implicitly represent its own policy --- a remarkable instance of a transformer encoding reflexive self-knowledge.

\paragraph{Cross-regime training dynamics.} The probes above can be applied not only to fully-trained models but to checkpoints throughout training, giving us the time-course of representation emergence in each regime. The hypothesis predicts a specific dissociation pattern. World-state belief should sharpen in all five regimes, since it is required for any observation prediction. The two notions of Self representation we distinguished above should dissociate across regimes. A \emph{token-type} Self, distinguishing action positions from observation positions for the purposes of belief updates, should emerge in Regimes 2--4, where the model produces its own actions; the action-position entropy collapse should be sharpest in (2a). An \emph{actor-identity} Self, distinguishing the model's own action stream from those of other actors, should emerge specifically in Regime 5, where the model is exposed to multiple actors and must distinguish among them to predict observations. A clean sub-question within token-type Self is the contrast between (2a) and (2b): if a Self representation emerges in (2b), where the loss never directly touches action positions, then the causal-sampling structure of the input is itself sufficient; if only in (2a), then explicit gradient at action positions is required. Either outcome refines what we should expect to find in RL-trained reasoning LLMs, where both pressures are typically present at once.

\paragraph{Self-recognition and off-policy generalisation.} Actor-identity self-recognition is not predicted in Regimes 2--4, where the model is always the actor and has no contrastive training signal to distinguish itself from another agent. Regime 5 changes this by exposing the model to action streams from many policies including its own, giving the actor-identity Self representation a functional purpose. The first test, then, is whether ``this trajectory was produced by my own policy'' is linearly decodable from the residual stream in Regime 5 but not in Regimes 2--4. A positive result in Regime 5 is the toy-model analogue of the on-policy detection finding in our forthcoming work with Lindsey, and is the expected outcome if the actor-identity Self has formed. The harder test is off-policy generalisation: when presented with an action stream from a policy not in the training family, does the model produce correct Bayes-adaptive predictions? This distinguishes two hypotheses about the kind of computation the model has learned. Under Hypothesis A, the model implements the abstract action-conditioned belief update, so the source of the action does not matter and generalisation follows. Under Hypothesis B, the model has learned a policy-specific shortcut through learning the joint statistics of (action, observation) under the policies it has seen and generalises poorly.

\paragraph{Causal control via Self-steering.} To establish that the SWITCH factorisation is causal, we perform targeted interventions on the Self and World subspaces identified above. Writing a different Self encoding into the residual stream should shift the action distribution in the direction the SWITCH dynamics predict, without disturbing the World belief; writing a perturbed World encoding should distort observation predictions and downstream action choices in the direction the Bayes-adaptive update predicts. A clean dissociation across these interventions is strong evidence that the SWITCH factorisation is real and not merely linearly decodable, and provides a toy-model existence proof that the Self component of an analogous LLM representation could be a steering handle for interpretability and alignment.

\paragraph{Cross-world generalisation.} If the Self and World components are genuinely factored, they should be modifiable independently. We test this by training the model in a meta-learning regime over a family of SWITCHes sharing the same World-state cardinality and observation alphabet, but with transition matrices sampled from a meta-distribution and presenting a held-out SWITCH from the same family at evaluation. The hypothesis predicts that the World representation adapts in-context to the new task while the Self representation remains stable. This is a toy analogue of learning the self as a fixed actor in a volatile world.

\section*{Extension to Language Models (Second Half of the Year)}

The second half of the grant is conditional on the first half delivering interpretable findings on the SWITCH. Assuming it does, we extend the probe suite to RL-trained reasoning language models, but in a setting that preserves the genuine self/world interaction structure the hypothesis depends on. Pure chain-of-thought reasoning is not the best setting for such an experiment: the model's own tokens are its only input, and there is no external World with hidden state and action-conditioned dynamics for the model to attribute to. We focus instead on \emph{code-execution-augmented reasoning}, where the model interleaves natural-language reasoning with code that is actually executed and whose output is fed back into the context. The Python interpreter plays the role of World, with its own hidden state and action-conditioned responses, and the model's code blocks play the role of Self-emitted actions. Several open-source checkpoints exist in this configuration with a paired base model available (e.g., from the rStar-Math / Qwen-Math-with-code family), and we can check for the emergence of self/world factorizations during post-training.

Taken together, the proposed project offers a theoretical framework in the SWITCH model for a Self/World factorisation in intelligent agents, a mechanistic hypothesis for how this factorisation emerges from training pressures, and a concrete experimental design to test the hypothesis in a toy model and then in language models. This builds on a conceptual distinction between passive prediction and active exploration that is central to agentic learning, and a recognition of the special role that exploration plays in shaping the internal representations of an agent, and especially its self model. The project also has implications for interpretability and alignment, since the self-representation we predict is a natural handle for understanding and influencing the agent's behaviour.

\bibliography{references}

\end{document}
